\chapter{Worms (Intestinal Parasites)}

\section*{Definition}
Intestinal worms are parasitic organisms that live in the human digestive tract and feed on nutrients from the food you eat. The most common types are pinworms (threadworms), roundworms, hookworms, tapeworms, and whipworms. While common in children, they can affect people of all ages.

\section*{What Happens in Your Body}
Intestinal worms enter the body through the mouth (ingesting eggs from contaminated food, water, or hands) or through the skin (hookworm larvae burrow through bare feet). Once inside the digestive tract, they attach to the intestinal wall and grow, feed, and reproduce. Female worms lay eggs that are passed out of the body in stool. The type and severity of symptoms depend on the type of worm, the number of worms, and how long the infection has been present. Pinworms cause intense anal itching, especially at night, when the female worm crawls out to lay eggs on the surrounding skin.

\section*{Causes}
\begin{itemize}
  \item Pinworms (Enterobius vermicularis): Most common in school-aged children; spread through contact with contaminated surfaces (toys, bedding, toilet seats) and then touching the mouth
  \item Roundworms (Ascaris lumbricoides): From ingesting eggs in contaminated soil, food, or water; common in areas with poor sanitation
  \item Hookworms (Ancylostoma duodenale, Necator americanus): Larvae penetrate the skin from contaminated soil (walking barefoot) and travel through the bloodstream to the intestines
  \item Tapeworms (Taenia species): From eating undercooked beef, pork, or fish that contains tapeworm larvae
  \item Whipworms (Trichuris trichiura): From ingesting eggs in contaminated soil or food
  \item Threadworms (Strongyloides stercoralis): Unusual because it can complete its life cycle inside the body and cause long-lasting infection
\end{itemize}

\section*{When to See a Doctor}
See your doctor if:
\begin{itemize}
  \item You or your child has anal itching, especially at night — the most common symptom of pinworms
  \item You see worms or worm segments in your stool or around your child's anus
  \item You have unexplained abdominal pain, nausea, or diarrhea
  \item You have weight loss despite normal appetite
  \item Your child is irritable, restless, or has trouble sleeping (common with pinworms)
  \item You have traveled to an area with known parasitic infections and develop digestive symptoms
  \item You have anal itching that affects the entire family — pinworm infections often spread within households
\end{itemize}

\section*{Related Symptoms}
\begin{itemize}
  \item Intense anal itching, especially at night (pinworms)
  \item Restless sleep or irritability
  \item Visible worms in stool (look like small white threads — pinworms; or long, round, spaghetti-like — roundworms; or flat, ribbon-like segments — tapeworms)
  \item Abdominal pain or cramping
  \item Nausea or vomiting
  \item Diarrhea or loose stools
  \item Weight loss
  \item Fatigue
  \item Anemia (from hookworms feeding on blood)
  \item Skin rash or itching at the site where hookworm larvae entered the skin
\end{itemize}
\textbf{Important:} The classic sign of pinworm infection is intense anal itching at night — if you suspect pinworms, check the anal area with a flashlight about 2-3 hours after the child falls asleep; you may see tiny white threads (the female worms).

\section*{Self-Care / Home Management}
\begin{itemize}
  \item Practice strict hand hygiene — wash hands thoroughly with soap and warm water after using the bathroom and before eating.
  \item Keep fingernails short and clean to reduce eggs trapped under nails.
  \item Wash all bedding, pajamas, towels, and underwear in hot water and dry on high heat.
  \item Vacuum carpets and clean bathroom surfaces thoroughly.
  \item Do not share towels or washcloths with an infected person.
  \item Shower in the morning to wash away eggs laid overnight.
  \item All household members should be treated if one person has pinworms — reinfection is very common without treating everyone.
  \item Refrigerate all family members' toothbrushes separately or soak them in hot water after use.
\end{itemize}

\section*{How Your Doctor Will Evaluate You}
\begin{itemize}
  \item For pinworms: The "tape test" — press a piece of clear tape against the anal area first thing in the morning (before showering or using the toilet), then bring the tape to your doctor to examine under a microscope for eggs.
  \item Your doctor may also examine the anal area with a flashlight to look for adult worms.
  \item A stool sample can be examined for eggs, larvae, or worm segments to identify other types of worms.
  \item Blood tests may show anemia (with hookworm infection), elevated eosinophils (a type of white blood cell that increases with parasitic infections), or specific antibodies.
  \item A complete history of travel, diet (especially undercooked meat or fish), and exposure to soil or children is important for diagnosis.
\end{itemize}

\section*{Treatment Options}
\begin{itemize}
  \item Anti-parasitic medications are very effective and usually require just one or two doses:
    \begin{itemize}
      \item Pinworms: Mebendazole or albendazole — one dose, repeated after two weeks
      \item Roundworms, hookworms, whipworms: Mebendazole, albendazole, or ivermectin
      \item Tapeworms: Praziquantel — one dose
      \item Threadworms: Ivermectin or albendazole (often requires longer treatment)
    \end{itemize}
  \item All household members should be treated for pinworms even if they have no symptoms — reinfection cycles easily.
  \item The itching from pinworms may persist for about a week after treatment even though the worms are gone; antihistamines or soothing creams can help.
  \item Re-treat after two weeks for pinworms to catch any newly hatched worms.
  \item Good hygiene after treatment is essential to prevent reinfection.
  \item Iron supplements may be needed for hookworm-related anemia.
  \item In severe cases with complications (bowel obstruction from heavy roundworm infection), hospitalization may be needed, but this is rare.
\end{itemize}

\section*{References}
\begin{thebibliography}{99}
\bibitem{worms:ref1} Bethony J, Brooker S, Albonico M, et al. "Soil-transmitted helminth infections: Ascariasis, trichuriasis, and hookworm." \textit{The Lancet}, 2006;367(9521):1521--1532.
\bibitem{worms:ref2} Jourdan PM, Lamberton PHL, Fenwick A, et al. "Soil-transmitted helminth infections." \textit{The Lancet}, 2018;391(10117):252--265.
\bibitem{worms:ref3} Ash LR, Orihel TC. \textit{Atlas of Human Parasitology}, 5th ed. American Society for Clinical Pathology Press, 2007.
\bibitem{worms:ref4} Kucik CJ, Martin GL, Sortor BV. "Common intestinal parasites." \textit{American Family Physician}, 2004;69(5):1161--1168.
\bibitem{worms:ref5} Else KJ, Keiser J, Holland CV, et al. "Whipworm and roundworm infections." \textit{Nature Reviews Disease Primers}, 2020;6(1):44.
\bibitem{worms:ref6} Keiser J, Utzinger J. "The drugs we have and the drugs we need against major helminth infections." \textit{Advances in Parasitology}, 2010;73:197--230.
\end{thebibliography}
